\section{Limitations and Conclusion}

The evidence covers one model, one H100, one fixed serving configuration,
fixed 128-token outputs, and workload transfer only.  It does not establish
configuration ranking, finite-rate queueing, quality equivalence, power-cap
transfer, TP/PP fidelity, H200/B200 transfer, or net deployment energy savings.
Six development workloads are small for a four-feature model.  The
preregistered low-concurrency TTFT failure shows that a simple range envelope
is too permissive near sparse corners, while the resulting $k<4$ abstention is
validated only for this workload family.  DCGM measures GPU energy but excludes host,
network, cooling, and embodied costs.  The current prediction intervals are
engineering bounds, not calibrated confidence intervals.  CPU projection
energy and queue-time savings are also unmeasured, so ``zero-GPU'' denotes
resource allocation, not net energy.

Within that boundary, \system{} reaches \CalibratedEnergyMAPE{} and
\ConfirmationEnergyMAPE{} energy MAPE on two sealed holdouts with no intervening
refit.  The second study also converts a preregistered latency failure into an
explicit abstention rather than a confident recommendation.  The central
result is not that a Docker container simulates an H100.  It is that CPU-first
screening can guide scarce GPU experiments only when prediction, provenance,
scope decisions, and real validation are one system.  The next step is to test
whether this contract reduces measured GPU-hours for SLO-constrained
configuration and topology search on HPC clusters.
